\section{Motivação e Oportunidades de Evolução do Projeto}
\label{sc:motivacao}

Os bons resultados obtidos com a plataforma na Iniciação Científica motivaram a continuidade da pesquisa neste TCC, buscando enriquecer o sistema por meio de quatro frentes de aprimoramento:

\textbf{1. De Painéis de Visualização para Alertas Automáticos:} A versão atual conta com painéis interativos detalhados no Apache Superset, nos quais o usuário explora mapas e filtros na tela. Adicionar um mecanismo que envie alertas automáticos assim que um evento atingir alto risco (sem depender que alguém esteja olhando o painel naquele instante) amplia a capacidade do sistema de apoiar brigadistas e equipes da Defesa Civil em tempo hábil \cite{endsley1995toward}.

\textbf{2. Consultas em Linguagem Natural e Relatórios Prontos:} Embora os painéis visuais cubram as visões principais, investigações específicas sobre o banco de dados hoje exigem escrever consultas técnicas em SQL. A inclusão de um assistente conectado ao banco permite que um profissional faça perguntas em português e receba imediatamente um boletim padronizado em PDF, facilitando o uso da ferramenta por pessoas fora da área de computação.

\textbf{3. Inclusão do Vento e da Vegetação na Direção do Fogo:} Na versão atual, o sistema mede o risco de proximidade usando um raio igual para todos os lados (um círculo de 5 km a 15 km ao redor do fogo). Porém, ventos fortes (que chegaram a 48,2 km/h em episódios recentes segundo o ERA5) e vegetações secas de pastagem fazem o fogo correr muito mais rápido a favor do vento \cite{mapbiomas2026fogo}. Considerar a direção do vento permite esticar a área de alerta para o lado correto e reduzir alarmes desnecessários no sentido contrário.

\textbf{4. Aprofundamento Analítico com o Histórico Oficial do Setor Elétrico:} Apesar de o projeto monitorar hospitais do SUS, rodovias e terras indígenas, o setor elétrico oferece uma oportunidade especial de validação estatística por contar com o acervo público do ONS, que registra 8.395 interrupções de carga entre 2007 e 2026 (717 GW acumulados) \cite{ons2026dadosabertos}. Cruzar os dados gerados pela nossa plataforma com esse histórico oficial permite demonstrar quantitativamente como os padrões de queimadas afetam a Rede Básica ($\ge 230$ kV).
